\color{unimain}{\subsection{Управление выключателем}}
\color{black}

\begin{enumerate}


\item
Функциональный блок управления выключателем предназначен для подачи сигнала на включение силового выключателя.  

\item
Запрет на включение БСК формируется:

\begin{itemize}
    \item После разряда БСК с уставкой выдержки времени \textbf{«Тожид. Разряда БСК»}. Выводится накладкой \textbf{«Ожид. Разр. БСК»}. 
    \item При срабатывании логики отключения основных защит;
    \item При срабатывании логики отключения резервных защит
    \item При срабатывании логики отключения от режимных защит.
\end{itemize}

\item
Команда на включение выключателя \textit{«Включить В1»} формируется от команды оперативного включения выключателя (сигнал \textit{«РКВ В1(ЛВ)ВН»}).

\medskip
\begin{figure}[H]
\centering
\includegraphics[width=0.9\textwidth,height=0.9\textheight,keepaspectratio]{\fbpath/09.01 БСК/761.3391 УВ/img/UV.pdf}
\captionof{figure}{Функциональная схема логики отключения выключателя}\label{pict_uv:UV}
\end{figure}

\item
Выдержка времени на возврат \textbf{«Тпрод опер. вкл. В1»} задает время продления сигнала оперативного включения. 

\item
При необходимости может быть введен подхват команды на включение от токового реле, установленного в цепи ЭМВ, в этом случае сигнал на включение удерживается в сработанном состоянии до исчезновения тока через ЭМВ (пока блок-контакты выключателя не разорвут цепь включения). Ввод осуществляется программным ключом \textbf{«Подхват вкл. В1 от РТ ЭМВ»}. 

\item
Предусмотрена возможность блокировки многократных включений на КЗ. Если при наличии сигнала включения будет протекать ток по электромагнитам отключения (что соответствует включению на КЗ), то сигнал \textit{«Включить В1»} заблокируется, пока не пропадут сигналы включения (от АПВ или оперативного включения). Блокировка вводится программным ключом \textbf{«Блок. от многокр. вкл. на КЗ В1»}. 

\item
При наличии сигнала \textit{«Блок. включения В1»}, формируемого функциональным блоком контроля силового выключателя, включение выключателя блокируется. 


\end{enumerate} 